\ourchapter{Background}
\label{chapter:background}

This chapter discusses preliminary concepts that are utilised throughout the remainder of the project, some of which already feature in available tools for reviewing code. 

\section{Static Code Analysis}

\textbf{Static code analysis} is the automated examination of source code without executing a program~\cite{StaticCodeAnalysis}, rather than \textbf{dynamic code analysis}, which examines the properties of a running program~\cite{DynamicCodeAnalysis}. In software development it is used primarily as a scalable method, performed early in a program's life cycle, to identify potential faults, unreachable code or security vulnerabilities by matching common patterns~\cite{SCAVulnerabilities}.

These static analysis tools will generate an abstract representation of a program, parsing through its code for syntactic and semantic information. This information is then reflected in a simplified model composed of aspects relevant to the tool's task, such as an \textbf{AST}.

\subsection{Abstract Syntax Trees (ASTs)}

An \textbf{abstract syntax tree (AST)} is a data structure, generated by a parser, that acts as a hierarchical representation of the abstract syntactic structure of source code~\cite{ASTs}. These trees (a type of graph as per \cref{section:graphs}) contain nodes that represent constructs like functions, variables and classes from within code. This allows them to be used in analysing features and structures within programs.

The AST in \Cref{fig:ast} below is an approximation of the result of parsing the Java method \emph{restoreFoo( )} in \Cref{lst:java_ast_example}:

\begin{figure}[H]
    \centering
    \includesvg[width=12cm]{background/AST.svg}
    \caption{Abstract Syntax Tree of the restoreFoo( ) method}
    \label{fig:ast}
\end{figure}

\begin{listing}[H]
\begin{minted}[linenos, breaklines]{java}
public void restoreFoo() {
    if (foo < 5) {
        foo = 10;
    }
}
\end{minted}
\caption{A simple Java function, \emph{restoreFoo}}
\label{lst:java_ast_example}
\end{listing}

An AST could be generated for student work as part of an automated tool, which would then be able to analyse the interactions between elements to detect more complicated, structural patterns.

\section{Object-Oriented Programming (OOP)} 
Created in the 1960s, the programming paradigm that is \textbf{Object-Oriented Programming (OOP)} revolves around the notion of objects, which encapsulate procedures and data by performing computations and saving a local state respectively~\cite{OOPThemes}. In contrast to traditional, procedural programming being a sequence of actions, OOP uses a collection of interdependent components that provide a service specified in their interfaces, which then relate to each other by client/server relations~\cite{OOPConcept}.

The goal of OOP is to structure programs in a way that more closely models real-world scenarios, which in turn make it easier for programmers to understand and develop. As a result, popular \textbf{object-oriented languages} like Java, Python and C\# are often used for software engineering courses, not just because OOP is prevalent within industry, but also because common problems can be abstracted as patterns that can be applied to multiple situations, and thus reused across programs.

\subsection{Design Patterns}

These \textbf{design patterns} were famously formalised by the `Gang of Four' as ``identifying classes, instances, their roles, collaborations, and the distribution of responsibilities''~\cite{GangOfFour}. They proposed 23 foundational design patterns, categorised into three main types: creational, structural, and behavioural, each of which addresses different aspects of software design. Creational patterns target object instantiation mechanisms, structural patterns focus on the composition of classes and other objects, and behavioural patterns are concerned with object interactions and responsibilities.

For example, the \textbf{Singleton} pattern is a creational pattern that ensures that a given class only has a single instance, whilst providing a global point of access to that instance, typically through a class method called \emph{getInstance( )}. The Singleton pattern will typically be used in cases where a single instance of a class is needed to control other classes across a program, or in the case that a particular class is expensive to instantiate. A basic implementation of this pattern is displayed in \Cref{lst:singleton-pattern}.

\begin{listing}[H]
\begin{minted}[linenos, breaklines]{java}
public class Singleton {
    static Singleton instance = new Singleton();

    Singleton() {}

    public static Singleton getInstance() {
        return instance;
    }
}
\end{minted}
\caption{A minimal Java implementation of the Singleton pattern}
\label{lst:singleton-pattern}
\end{listing}

These patterns are essential knowledge for developers, as their use facilitates modular, reusable, and adaptable code. As such their application and purpose is commonly taught as part of higher education courses focused on software development, where one common method of representing these patterns is through UML diagrams~\cite{GraphicalDepictionOfDP}.

\subsection{Unified Modelling Language (UML) Diagrams}
Unified Modelling Language (UML) diagrams are a standardised format for the graphical representation of software design and structure. UML diagrams are particularly useful for showing design patterns as they can succinctly depict the relationship between classes and their components. Below is a basic UML diagram representing the use of the Java Singleton class in \Cref{lst:singleton-pattern}. For reference, - and + denote private and public visibility respectively, with methods separated from fields by use of brackets ( ).

\begin{figure}[H]
    \centering
    \includesvg[width=12cm]{background/SingletonUML.svg}
    \caption{UML diagram representing the Singleton pattern}
    \label{fig:uml}
\end{figure}

Design pattern UML diagrams provide an easily composed and understood representation of potentially complex relations. Whilst it is not entirely formalised, and as such still open to interpretation, it is possible to automatically recognise patterns from a diagram~\cite{UMLDesignPatterns}. This suggests that for an automated review tool, structural patterns could be similarly composed as part of a task specification and then detected within programs.

\section{Graph Theory}
\label{section:graphs}
\textbf{Graphs} provide a formal representation for networks. A graph, G, can be represented by the pair (N, E) where N is a set of \textbf{nodes}, representing objects and E is a set of \textbf{edges}, representing relations between pairs of objects, or more formally E $\subseteq$ \{\{m, n\} $\mid$ m, n $\in$ N\}.

When N = \{A, B, C, D, E\} and E = \{(A, B), (A, C), (A, D), (A, E), (C, D), (D, E)\}, the resulting graph, G, would be:

\begin{figure}[H]
    \centering
    \includesvg[width=8cm]{background/Graph.svg}
    \caption{An example of a simple graph, G}
    \label{fig:example-graph}
\end{figure}

The \textbf{degree} of a node within a graph is the number of edges incident on it, meaning the number of edges that have that node as an endpoint. In a \textbf{directed graph}, where edges are ordered relations, we then have the notion of \textbf{in} and \textbf{out} degrees, where we calculate the number of edges that end and begin at a node respectively. 

A \textbf{tree} is a specialisation of a directed graph, meaning we can apply some elements of graph theory to ASTs representing source code.

\subsection{Subgraphs}
For graphs G$_1$ and G$_2$, G$_1$ is a \textbf{subgraph} of G$_2$, or G$_1$ $\subseteq$ G$_2$, if nodes(G$_1$) $\subseteq$ nodes(G$_2$) and edges(G$_1$) $\subseteq$ edges(G$_2$). That is, if all nodes and edges within G$_1$ are also contained within G$_2$. 

For a graph, G', with N' = \{A, C, D, E\} and E' = \{(A, C), (A, D), (D, E)\}, we have that G' $\subseteq$ G:

\begin{figure}[H]
    \centering
    \includesvg[width=15.5cm]{background/SubGraph.svg}
    \caption{Graph G' (left) illustrated as a subgraph of G (right)}
    \label{fig:example-subgraph}
\end{figure}

\subsection{Graph Isomorphism}
An \textbf{isomorphism} is a bijective homomorphism~\cite{LatticeTheory}, meaning that for G to be isomorphic to a graph, G'', there must be a function f : nodes(G) $\rightarrow$ nodes(G'') and a function g : edges(G) $\rightarrow$ edges(G''), such that for any edge e $\in$ edges(G) with endpoints n$_1$ and n$_2$, then for g(e) the endpoints must be f(n$_1$) and f(n$_2$). 

For G'', with N'' = \{1, 2, 3, 4, 5\} and E'' = \{(1, 2), (1, 3), (1, 4), (1, 5), (3, 4), (4, 5)\}, we have that G'' is \textbf{isomorphic} to G in \Cref{fig:example-graph} (and the reverse, implicitly), as the nodes have simply been transformed to their respective position in the alphabet (A $\rightarrow$ 1, B $\rightarrow$ 2...) but the relations between them have remained the same (even if they might look different).

\begin{figure}[H]
    \centering
    \includesvg[width=7.5cm]{background/Isomorphism.svg}
    \caption{Graph G'', an isomorphism of G}
    \label{fig:example-isomorphism}
\end{figure}

These two concepts are used within \textbf{graph mining}, where abstract patterns can be detected within a larger structure by finding a pattern isomorphism in a subgraph~\cite{GraphMining}. As a result, this technique could be used to find instances of a given code structure within an AST. 

\section{Large Language Models (LLMs)}
\label{section:LLMs}
A \textbf{large language model (LLM)} is a machine learning model that can understand and produce natural language, approximating a human's response to queries. LLMs are trained for a considerable length of time over large datasets that span a wide range of text sources, such as books, articles and internet forums. This breadth of information allows LLMs to develop a broad understanding of language and context for a variety of subjects, to the point where they can approximate human experts. For instance, OpenAI's GPT-4, released in 2023, is capable of ``passing a simulated bar exam with a score around the top 10\% worldwide''~\cite{GPT4Report}. As ChatGPT was trained using data from Stack Overflow among other online resources, it is relatively adept at understanding and critiquing source code, making it potentially applicable in an automated feedback tool.

However, beyond the issue of the computational power required for training, with its corresponding time requirements and substantial energy consumption, LLMs suffer from the potential for bias due to the data used in training and `hallucinations', where ``models generate plausible-sounding but unfaithful or nonsensical information''~\cite{LLMHallucinations}. Such hallucinations can cause issues in applications where correctness is critical, such as the medical and legal fields. There are methods to combat these issues, such as providing an LLM with more useful input through \textbf{prompt engineering}.


\subsection{Prompt Engineering}
Prompt engineering is the discipline of ``optimizing prompts to effectively utilize LLMs''~\cite{PromptEngineering}, a prompt being the natural language request given to an LLM. This can be accomplished by providing the LLM a tailored frame of reference, increased specificity in requests, supplying additional context and examples, and clearly specifying the format a desired output should take. One specific approach is \textbf{few-shot learning}, where LLMs are provided examples to learn from as part of a prompt, which normally provides more positive results~\cite{FewShot}, but requires that examples are available and accurate.

\section{Educational Technology (EdTech)}

\textbf{Educational Technology}, or \textbf{EdTech}, is the use of hardware and software to facilitate learning by improving accessibility, engaging students and allowing for a wide range of assessments~\cite{20YearsEdTech}.

With the popularisation of machine learning, \textbf{adaptive learning platforms} have become a growing form of EdTech, constantly evaluating learners' performance to produce individualised learning paths~\cite{AdaptiveLearningPlatforms}. This tailored approach to learning could be incorporated into automated code analysis by providing users with feedback relevant to their code's current state, rather than overwhelming them with an exhaustive list of errors.
